%%%%%%%%%%%%%%%%%%%%%%
\documentclass{singlecol-new}
\usepackage[utf8]{inputenc}
\usepackage{grffile}


\usepackage{natbib,stfloats}
\usepackage{mathrsfs}

\def\newblock{\hskip .11em plus .33em minus .07em}

\theoremstyle{TH}{
\newtheorem{lemma}{Lemma}
\newtheorem{theorem}[lemma]{Theorem}
\newtheorem{corrolary}[lemma]{Corrolary}
\newtheorem{conjecture}[lemma]{Conjecture}
\newtheorem{proposition}[lemma]{Proposition}
\newtheorem{claim}[lemma]{Claim}
\newtheorem{stheorem}[lemma]{Wrong Theorem}
\newtheorem{algorithm}{Algorithm}
}

\theoremstyle{THrm}{
\newtheorem{definition}{Definition}[section]
\newtheorem{question}{Question}[section]
\newtheorem{remark}{Remark}
\newtheorem{scheme}{Scheme}
}

\theoremstyle{THhit}{
\newtheorem{case}{Case}[section]
}

\makeatletter
\def\theequation{\arabic{equation}}

%\JOURNALNAME{\TEN{\it Int. J. System Control and Information
%Processing,
%Vol. \theVOL, No. \theISSUE, \thePUBYEAR\hfill\thepage}}%
%
%\def\BottomCatch{%
%\vskip -10pt
%\thispagestyle{empty}%
%\begin{table}[b]%
%\NINE\begin{tabular*}{\textwidth}{@{\extracolsep{\fill}}lcr@{}}%
%\\[-12pt]
%Copyright \copyright\ 2012 Inderscience Enterprises Ltd. & &%
%\end{tabular*}%
%\vskip -30pt%
%%%\vskip -35pt%
%\end{table}%
%}
\makeatother

%%%%%%%%%%%%%%%%%
\begin{document}%
%%%%%%%%%%%%%%%%%

\setcounter{page}{1}

%\LRH{A. Pardeshi et~al.}

\RRH{From Signal Fusion to Asset Allocation}

\VOL{x}

\ISSUE{x}

\PUBYEAR{201X}

\BottomCatch

%\CLline

\subtitle{}

\title{From Signal Fusion to Asset Allocation: A Decision-Theoretic Framework for Regime-Conditional Portfolio Construction}

%
%\authorA{Anandkumar Pardeshi}
%
%affA{Department of Information Technology,\\ Fr. C. Rodrigues Institute of Technology,\\ Vashi, Navi Mumbai \\
%E-mail: anand.pardeshi@fcrit.ac.in}
%
%
%\authorB{Sujata Deshmukh}
%
%\affB{Department of Computer Engineering,\\ Fr. C. Rodrigues College of Engineering,\\ Bandra, Mumbai \\
%E-mail: sujata.deshmukh@fragnel.edu.in}
%

%
%\authorA{Fusheng Wang\footnote{Work done while working at Siemens Corporate Research.} }
%
%\affA{Department of Biomedical Informatics, Emory University
%\newline
%36 Eagle Row, Ste 589, Atlanta, GA 30322, USA}
%
%
%
%\authorB{\footnotesize Cristobal Vergara-Niedermayr\footnote{Work done while working at Siemens Corporate Research.}}
%\affB{Oracle \newline
 %New Jersey, USA}
%
%


\begin{abstract}
Predictive modelling in finance has advanced faster than the machinery that converts a
prediction into a position. This paper specifies and tests a decision-theoretic
framework for that conversion. Latent market states are identified by a Gaussian
mixture model fitted on an expanding window; technical and market-state composites are
fused with state-dependent weights and scaled to a return metric; the fused signal
enters a constrained quadratic program with covariance shrinkage and an explicit
turnover penalty; and the resulting portfolio is scaled by a regime-dependent risk
budget. Evaluated on ten liquid multi-asset instruments over sixteen years under a
walk-forward protocol with transaction costs, the framework attains a Sharpe ratio of
0.74 against 0.53 for an equal-weighted benchmark and reduces maximum drawdown from
$-36.1$ to $-19.0$ per cent, at the cost of a lower compound return. It outperforms in
all four stress episodes in the sample and in 18 of 18 specifications examined, though
the Sharpe improvement is not statistically significant. Ablation at matched average
exposure isolates the source: timing risk by market state improves both the Sharpe
ratio and the drawdown with no change in mean return, whereas the signal fusion layer
contributes nothing and the state-dependent weighting hypothesis is rejected in the
opposite direction. The paper also identifies two failure modes in frameworks of this
kind: an absolute volatility target that never binds, and a signal-to-allocation
conversion in which mismatched units rather than information determine the solution.
\end{abstract}

\KEYWORD{Portfolio Optimisation; Regime Detection; Risk Budgeting;
Decision-Theoretic Framework; Signal Fusion; Ablation Analysis; Asset Allocation}

\REF{}

\begin{bio}

\end{bio}


\maketitle

\section{Introduction}
\label{sec:intro}

Quantitative finance has invested heavily in prediction and comparatively little in the
step that follows it. A large body of work now forecasts returns, directions or
volatilities with machine learning, and reports its success in statistical terms:
accuracy, error, or the significance of a coefficient \citep{gu20, zhang22}. Far less
work specifies what an investor should then hold. That second step is not a
formality. A forecast becomes an investment only after it has been reconciled with a
covariance structure, a risk budget, position limits and the cost of trading, and the
translation can destroy the value of a forecast as easily as preserve it. This paper
refers to the missing step as the \emph{signal--allocation gap}.

\subsection{The Signal--Allocation Gap}

The gap is visible in how results are reported. Systematic reviews of fusion-based
prediction find that statistical accuracy remains the dominant success criterion, with
economic value reported inconsistently or not at all \citep{zhang22}. Where allocation
is addressed, it is often through heuristics rather than optimisation: the IntPort
framework of \citet{sami25}, for instance, converts forecasts into weights by
statistical averaging, without a covariance structure or an explicit risk objective.
Reinforcement learning approaches optimise a policy directly and so avoid the
translation problem, but the majority are formulated for a single instrument
\citep{saijai26, jiang25}, which leaves the portfolio-level question open.

This matters because the properties an investor cares about are portfolio properties.
Diversification, drawdown and the stability of holdings are not visible at the level of
an individual forecast, and \citet{demiguel09} showed that optimisation which ignores
estimation error can underperform naive equal weighting out of sample. Any framework
that claims to bridge prediction and allocation must therefore be judged on portfolio
outcomes under realistic frictions, not on the quality of its forecasts.

\subsection{Regime Dependence}

Financial markets do not behave homogeneously through time. Since \citet{hamilton89},
regime-switching models have been used to represent discrete shifts in the statistical
behaviour of financial series, and \citet{angbekaert02} demonstrated that accounting for
such shifts materially changes optimal international allocations. \citet{kritzman12}
made the practical case that regime awareness alters what a dynamic strategy should do,
not merely how it should be described. More recent work continues in this direction,
combining regime identification with allocation rules \citep{li25, luo26}.

Two consequences follow for a framework of the kind proposed here. The first is that a
predictive signal's reliability is conditional: an indicator that is informative in a
quiet trending market need not be informative when the market is disturbed. The second
is that the appropriate amount of risk to carry is itself state-dependent, which is the
premise behind volatility-managed strategies \citep{moreira17, hood25} and behind risk
budgeting more generally \citep{asimit25}. A framework that detects regimes but acts on
them in only one of these two ways will leave part of its own premise untested.

\subsection{What This Paper Does}

This paper specifies a three-stage decision-theoretic framework and tests it end to
end. Latent market states are identified by a Gaussian mixture model fitted on an
expanding window. Two composite signals---one technical, one derived from price and
volume dynamics---are fused with regime-conditional weights and converted to expected
excess returns using information-ratio scaling. The fused signal then enters a
constrained quadratic program with Ledoit--Wolf covariance shrinkage
\citep{ledoit04} and an explicit turnover penalty, and the resulting sleeve is scaled by
a regime-dependent risk budget.

The framework is evaluated on ten liquid exchange-traded funds spanning equities, fixed
income, credit, gold and real estate over sixteen years, under a walk-forward protocol
with transaction costs. The evaluation is designed to answer not only whether the
framework works but which of its parts is responsible, and it reports the answer in
both directions. The regime risk budget is shown to be the effective component: at
identical average exposure, timing the de-risking by regime improves both the Sharpe
ratio and the maximum drawdown. The signal fusion layer is shown not to be: the
information coefficients of both composites are statistically indistinguishable from
zero or significantly negative, and the regime-conditional weighting hypothesis is
rejected by the data that motivated it.

Reporting the second finding is a deliberate choice. Component-level nulls are rarely
published, which is one reason architectures accumulate layers whose contribution has
never been isolated. \citet{harvey15} and \citet{valls26} document how easily an
unreported specification search inflates apparent performance in this setting, and
\citet{demiguel09} showed the practical cost of trusting an unvalidated component. An
ablation that locates the source of a result is more useful to a practitioner than an
aggregate number that does not.

\subsection{Contributions}

\begin{itemize}

\item A fully specified regime-conditional allocation framework, given in enough detail
to be reproduced from the text alone: exact feature definitions, the mixture model and
its label-anchoring requirement, the optimisation problem with all constraints, and the
walk-forward protocol.

\item A two-step separation of the allocation decision into relative composition and
total risk. The paper shows that expressing the regime risk budget as a fraction of the
portfolio's own volatility, rather than as an absolute volatility target, is what keeps
the regime layer active: absolute targets calibrated for concentrated equity portfolios
bind in one rebalance out of 192 on a diversified universe and render the layer inert.

\item An ablation design that isolates component contributions at matched risk. Because
scaling a portfolio toward cash changes both its return and its volatility, comparing a
de-risked strategy with a fully invested benchmark confounds risk timing with risk
level. Holding average exposure fixed separates the two, and shows that the regime
layer earns its place through timing rather than through simply holding less.

\item A reported null on regime-conditional signal fusion, together with the
information coefficients that produce it, measured overall and within each state.

\item A documented diagnosis of a failure mode in signal-to-allocation conversion:
passing an unscaled indicator into a mean-variance objective allows the ratio of units
between the linear and quadratic terms, rather than the information in the signal, to
determine the solution, producing a degenerate single-asset portfolio.

\end{itemize}

\subsection{Organisation}

Section~2 reviews portfolio construction, regime-aware allocation and the treatment of
textual sentiment, and positions this study against them. Section~3 specifies the
framework. Section~4 reports the results, including robustness, ablation and the
negative findings. Section~5 discusses implications and limitations, and Section~6
concludes.

\section{Literature Review}
\label{sec:lit}

\subsection{From Mean--Variance Optimisation to Risk-Based Allocation}

Modern portfolio construction begins with \citet{markowitz52}, who formalised
allocation as a trade-off between expected return and variance. The framework is
theoretically complete and practically fragile: its solutions are highly sensitive to
the expected-return vector, and sample covariance matrices estimated from short windows
on many assets are poorly conditioned in exactly the directions the optimiser exploits.
\citet{demiguel09} quantified the consequence, finding that a naive equal-weighted
portfolio outperformed a range of optimised alternatives out of sample once estimation
error was taken into account.

Two responses dominate. The first improves the inputs. \citet{ledoit04} introduced a
shrinkage estimator that trades a small bias for a large variance reduction and yields
a well-conditioned covariance matrix; the approach remains standard, and recent work
extends shrinkage to higher moments \citep{lu25}. The second abandons expected returns
altogether in favour of risk-based construction. Inverse-volatility weighting, risk
parity and hierarchical risk parity \citep{prado16} allocate on the basis of the risk
model alone, and \citet{asimit25} develop the risk-budgeting formulation further by
linking portfolio selection to risk sharing; \citet{khodamoradi23} pursue a related
robustness objective through multi-interval mean-conditional-value-at-risk
optimisation, and \citet{kehinde25} evaluate volatility targeting through an inverse
data-envelopment formulation. Volatility management is a related idea
applied at the portfolio level rather than the asset level: \citet{moreira17} showed
that scaling exposure inversely to recent volatility improves risk-adjusted returns
across a range of factors, and \citet{hood25} examine how much of that improvement is
attributable to trend following rather than to volatility management itself.

Table~\ref{tab:paradigms} positions the present study against these paradigms.
This study sits between the two responses. It retains an expected-return term, but
scales it explicitly so that its magnitude is commensurate with the risk term, and it
separates the composition decision from the total-risk decision so that the two can be
evaluated independently.

\begin{table}[!ht]
\centering
\caption{Portfolio construction paradigms and the position of this study}
\label{tab:paradigms}
\small
\begin{tabular}{p{2.4cm}p{2.4cm}p{2.6cm}p{2.6cm}}
\hline
\textbf{Paradigm} & \textbf{Representative work} & \textbf{Basis of allocation} & \textbf{Principal limitation} \\
\hline
Mean--variance optimisation & \citet{markowitz52} & Expected returns and covariance & Extreme sensitivity to estimated inputs \\
\hline
Naive diversification & \citet{demiguel09} & None; equal weights & Ignores risk structure entirely \\
\hline
Shrinkage-based optimisation & \citet{ledoit04, lu25} & Regularised covariance & Still requires expected returns \\
\hline
Risk-based construction & \citet{prado16, asimit25} & Risk model only & Forgoes signal information by design \\
\hline
Volatility management & \citet{moreira17, hood25} & Total risk scaled by recent volatility & Does not address composition \\
\hline
Regime-switching allocation & \citet{angbekaert02, kritzman12, luo26} & State-dependent parameters & Component contributions rarely isolated \\
\hline
Learning-based allocation & \citet{jiang25, ashrafzadeh25, agal25} & Learned policy or predicted returns & Interpretability; frequently single-asset \\
\hline
\textbf{This study} & --- & Regime-conditional signal, plus a relative regime risk budget & Signal layer tested and found not to contribute \\
\hline
\end{tabular}
\end{table}

\subsection{Regime Dependence and Regime-Aware Allocation}

The premise that financial series alternate between statistically distinct states dates
to \citet{hamilton89}, whose Markov-switching formulation remains the reference point.
\citet{angbekaert02} applied regime switching to international allocation and found that
correlations rise in bear states, weakening diversification precisely when it is most
needed---an observation that motivates state-dependent covariance and state-dependent
risk budgets alike. \citet{kritzman12} argued that regime awareness should change the
strategy rather than only the commentary, and \citet{barroso15} demonstrated a
closely related mechanism at the factor level, showing that managing the volatility of
a momentum strategy removes most of its crash risk.

Contemporary work continues along two lines. One identifies regimes with increasingly
flexible statistical machinery: \citet{li25} combine a jump model with allocation
rules, and \citet{luo26} construct dual-regime signals with regime-dependent
parameters. The other embeds state-dependence inside a learning system, where the
regime is implicit in the learned policy \citep{jiang25}. Across both lines, the common
gap is evidential rather than methodological. Regime-aware systems are usually
evaluated against regime-unaware benchmarks in aggregate, so the reader learns that the
system as a whole performs well but not which of its state-dependent components is
responsible. The ablation design in Section~4 is intended to close that gap for the
framework proposed here.

\subsection{Machine Learning and Multimodal Signals}

Machine learning entered this field early and has since become routine. Agent-based
and reinforcement formulations were applied to daily equity trading by \citet{lee07},
metaheuristic-optimised regression to price forecasting by \citet{chou18}, and pattern
mining with fuzzy inference to trading-point detection by \citet{huang20}; each
improved predictive accuracy without addressing what the resulting forecasts imply for a
portfolio. \citet{gu20} provide the benchmark comparative study in asset pricing,
finding consistent gains from flexible models while noting that measured predictability
is small. Applications to allocation
are more recent: \citet{ashrafzadeh25} combine deep learning with stock clustering for
portfolio construction, \citet{agal25} apply learned risk estimates to asset
allocation, \citet{ha26} compare machine-learning portfolios against a fixed
all-weather allocation, and \citet{singsiri26} build portfolios directly on money-flow
indicators of the kind used in this study. Reinforcement learning has been extended from single assets to
multi-period, high-dimensional allocation \citep{jiang25, saijai26}, and
\citet{wang26} couple prediction with an explicit trading strategy.

Fusion of heterogeneous signals is a parallel development. \citet{zhang22} review
decision fusion for equity prediction across a large body of primary studies and find
that fused architectures usually outperform unimodal ones, while also documenting wide
variation in fusion protocol and evaluation. \citet{farimani24} propose an adaptive
multimodal model that weights sources dynamically, and \citet{zong25} use gated
cross-attention to stabilise multimodal fusion for movement prediction. What these
studies establish is that fusion can help prediction. Whether a particular fusion rule
helps a particular allocation decision is a separate question, and one that requires a
component-level test rather than an aggregate comparison.

\subsection{Textual Sentiment and Market-Derived Proxies}

Textual sentiment entered finance through dictionary methods, and
\citet{loughran11} demonstrated that general-purpose sentiment lexicons misclassify
financial language badly enough to require domain-specific alternatives. Transformer
models and large language models have since displaced lexicons: \citet{officioso26}
combine language-model sentiment with market data, \citet{muhammad26} benchmark large
language models for target-based financial sentiment, \citet{mun25} derive investment
signals from language-model sentiment, and \citet{chuang25} compare language models on
financial news for sentiment and price prediction. Collectively these studies indicate
that genuine textual sentiment carries information beyond price history, and
\citet{mazzotta25} show more broadly that alternative data can carry information that
conventional variables do not.

They also indicate what such a study requires: a time-stamped corpus with coverage of
the traded instruments. That requirement is binding. For broad exchange-traded funds
there is no instrument-level news corpus of comparable coverage, which is why studies
of this kind are almost always conducted on individual equities. The present study
therefore uses market-derived proxies for attention and stress rather than textual
sentiment, states this explicitly in Section~3.3, and---rather than assuming the proxies
are adequate---measures their predictive content directly and reports that they are
not.

\subsection{Evaluation Practice}

A final strand concerns whether reported results should be believed.
\citet{harvey15} show that conventional significance thresholds are inadequate once the
number of strategies examined is taken into account, and \citet{valls26} document the
same problem specifically for volatility-based strategies, while \citet{sheppert26}
proposes an objective function designed to resist it. \citet{zhangz25} make the related
point that the model selected as optimal in sample is frequently not the one that
performs best out of it. \citet{ledoit08} provide
robust inference for Sharpe ratio comparisons under non-normal, autocorrelated returns,
and \citet{politis94} supply the stationary block bootstrap used here to construct
interval estimates that respect serial dependence. \citet{newey87} remains the standard
correction for autocorrelated and heteroskedastic test statistics.
\citet{raimundo26} survey the broader intersection of forecasting and portfolio theory
and reach a similar conclusion about the weight placed on point estimates.

This study adopts these practices as design constraints rather than as afterthoughts:
every specification examined is reported, inference is corrected for serial dependence,
interval estimates accompany point estimates, and the comparison against a costless
daily-rebalanced equal-weighted benchmark is deliberately unfavourable to the proposed
framework.

\section{Proposed Methodology}
\label{sec:method}

This section specifies the framework in full: the data, the exact form of every input
feature, the regime model, the fusion rule, the optimisation problem, and the
walk-forward protocol under which it is evaluated. The level of detail is deliberate,
so that the study can be reproduced independently from this description alone.

\begin{figure*}[t]
\centering
\includegraphics[width=\textwidth]{images/fig_architecture.png}
\caption{The proposed framework. Stage~1 constructs causal, cross-sectionally
standardised features; stage~2 identifies the market state and fuses the two composites;
stage~3 separates relative composition from total risk; stage~4 evaluates the result and
ablates its components.}
\label{fig:architecture}
\end{figure*}

\subsection{Data and Study Universe}

The primary study universe consists of ten exchange-traded funds selected to span the
major liquid asset classes available to a diversified allocator: US large capitalisation
equity (SPY), US large capitalisation growth equity (QQQ), US small capitalisation
equity (IWM), developed international equity (EFA), emerging market equity (EEM),
long-maturity US Treasuries (TLT), intermediate-maturity US Treasuries (IEF),
investment-grade corporate credit (LQD), gold (GLD) and US real estate (VNQ). Breadth
matters for the research question. A framework whose stated purpose is to redirect risk
across market states can only be tested where there is somewhere else for risk to go;
a single-sector universe offers no such destination, a point returned to in
Section~\ref{sec:megacap}.

Daily open, high, low, close and volume series are obtained from Yahoo Finance through
the \texttt{yfinance} interface, adjusted for splits and dividends. The download window
is 1~January 2007 to 1~January 2024. Indicator warm-up and the 252-session covariance
estimation window consume the first year, so the evaluated out-of-sample period runs
from January 2008 to December 2023 and covers the global financial crisis, the 2011
sovereign debt episode, the 2020 pandemic crash and the 2022 inflation-driven bear
market. A secondary universe of ten mega-capitalisation US equities over 2018--2024 is
carried through the same pipeline and reported separately in
Section~\ref{sec:megacap}.

\subsection{Feature Construction}

Nine features are computed for each asset on each session. All are causal: every value
at session $t$ uses only information available at the close of session $t$. Table~\ref{tab:features}
gives the exact definition of each, and no other input enters the framework.

\begin{table}[!ht]
\centering
\caption{Input features. $C_t$, $H_t$, $L_t$ and $V_t$ denote close, high, low and
volume. All features are computed per asset, then cross-sectionally standardised.}
\label{tab:features}
\small
\begin{tabular}{p{2.0cm}p{1.8cm}p{5.2cm}p{1.5cm}}
\hline
\textbf{Feature} & \textbf{Composite} & \textbf{Definition} & \textbf{Window} \\
\hline
RSI & Technical & Relative strength index, centred: $\mathrm{RSI}_t - 50$ & 14 \\
\hline
MACD & Technical & Moving average convergence divergence line, scaled by price: $\mathrm{MACD}_t / C_t$ & 12, 26, 9 \\
\hline
Stochastic & Technical & Stochastic oscillator \%K, centred: $\%K_t - 50$ & 14, smooth 3 \\
\hline
MFI & Technical & Money flow index, centred: $\mathrm{MFI}_t - 50$ & 14 \\
\hline
OBV flow & Technical & Twenty-session change in on-balance volume, normalised by average volume: $(\mathrm{OBV}_t - \mathrm{OBV}_{t-20}) / \overline{V}_{t,20}$ & 20 \\
\hline
Return & Market state & Five-session simple return: $C_t / C_{t-5} - 1$ & 5 \\
\hline
Relative volume & Market state & $V_t / \overline{V}_{t,20}$ & 20 \\
\hline
Realised volatility & Market state & Annualised standard deviation of daily returns, entered with a negative sign & 20 \\
\hline
Volume instability & Market state & Standard deviation of daily volume changes, entered with a negative sign & 20 \\
\hline
\end{tabular}
\end{table}

Each feature is standardised cross-sectionally on every session, so that a value
expresses how an asset ranks against its peers that day rather than against its own
history. Writing $x_{i,t}$ for the raw feature of asset $i$ and $N$ for the number of
assets,
\begin{equation}
z_{i,t} = \frac{x_{i,t} - \bar{x}_t}{s_t}, \qquad
\bar{x}_t = \frac{1}{N}\sum_{j} x_{j,t}, \qquad
s_t = \mathrm{sd}_j(x_{j,t}).
\end{equation}
Cross-sectional standardisation is used rather than time-series standardisation because
the allocation problem is a relative one: the optimiser must decide which assets to
prefer today, not whether today is unusual by historical standards. It also removes the
common market factor from the scores, so that a market-wide move does not appear as a
signal on every asset simultaneously.

The standardised features are then averaged into two composites. The technical
composite is the equally weighted mean of the five technical features. The market-state
composite is the equally weighted mean of the four market-state features, taken with
the signs given in Table~\ref{tab:features}, so that recent strength and active participation raise the
score while realised volatility and unstable volume lower it.

\subsection{What the Second Composite Is, and What It Is Not}
\label{sec:notsentiment}

The second composite is constructed entirely from price and volume. No news article,
social media post, corporate announcement or any other text is read at any point in
this study, and no lexicon, classifier or language model is applied. The composite is a
market-derived proxy for the same latent quantities that textual sentiment measures aim
to capture---attention, participation and stress---and it is reported here as exactly
that.

This is stated plainly because the distinction is material to how the results should be
read. Proxy construction of this kind is common in the literature and inexpensive to
compute, but it cannot be credited with the contextual information a transformer model
extracts from financial text, and it should not be described as though it could.
Section~\ref{sec:ic} measures the predictive content of both composites directly and
reports the outcome, including where that outcome is unfavourable. Extending the
framework to genuine textual sentiment requires a time-stamped news corpus with
coverage of the traded universe; for broad exchange-traded funds no such corpus exists
at instrument level, which is a limitation of the design rather than an oversight, and
Section~5 sets out what would be needed to remove it.

\subsection{Regime Identification}

Market states are identified with a Gaussian mixture model over two features, each
computed as a cross-sectional average across the universe: the twenty-session rolling
standard deviation of daily returns, and the ten-session rolling mean of daily returns.
The first captures the volatility of the environment, the second its direction. Both
are standardised before fitting.

The mixture uses $K = 3$ components with full covariance matrices, three random
restarts and a fixed seed. It is refitted at every rebalance date on an expanding
window containing only sessions strictly earlier than that date, with a minimum of 252
observations; the state of the current session is then obtained by prediction from that
fit. No observation ever contributes to the model that labels it.

One implementation detail is essential and is easy to overlook. The component indices
returned by a mixture fit are arbitrary and permute freely between refits, so labels
from consecutive fits are not comparable unless they are anchored. Components are
therefore sorted in ascending order of fitted mean volatility before use, which makes
regime~0 the calm state, regime~1 the transitional state and regime~2 the stress state
consistently across every refit. Omitting this step does not produce an error; it
produces a silently scrambled regime series and, if the resulting labels are compared
across fits, a spurious appearance of instability.

The choice $K = 3$ is not selected by information criterion. As
Section~\ref{sec:regimefit} reports, both the Akaike information criterion (AIC)
and the Bayesian information criterion (BIC) continue to fall as components are added
over the range tested, which is the usual behaviour of mixture models on
persistent financial series and reflects their tendency to absorb heavy tails with
additional components. The choice is instead made on persistence and interpretability:
three components yield states that last long enough to act on and that admit a clear
economic reading, whereas finer partitions fragment into short-lived states that a
monthly rebalancing schedule cannot exploit. This is a modelling judgement and is
reported as such.

\subsection{Regime-Conditional Signal Fusion}

The two composites are combined into a single score using weights that depend on the
prevailing regime, together with a per-asset confidence term. Let $T_{i,t}$ and
$M_{i,t}$ denote the standardised technical and market-state composites for asset $i$,
and let $r_t \in \{0,1,2\}$ be the regime. The fused score is
\begin{equation}
z_{i,t} = c_{i,t} \left[ w^{T}_{r_t} T_{i,t} + w^{M}_{r_t} M_{i,t} \right],
\end{equation}
with regime weights $(w^T, w^M)$ of $(0.7, 0.3)$ in the calm state, $(0.5, 0.5)$ in the
transitional state and $(0.3, 0.7)$ in the stress state. The direction of this schedule
follows the hypothesis stated in Section~1, that trend-following information is more
reliable in quiet trending markets while attention and stress information becomes
relatively more informative when markets are disturbed. Section~\ref{sec:ic} tests that
hypothesis rather than assuming it.

The confidence term $c_{i,t}$ down-weights assets whose two composites disagree:
\begin{equation}
c_{i,t} = \mathrm{clip}\!\left(0.5 + 0.5\,\rho^{(63)}_{i,t}, \; 0.1, \; 1.0\right),
\end{equation}
where $\rho^{(63)}_{i,t}$ is the trailing 63-session correlation between the two
composites for that asset. Agreement between independent views is treated as evidence;
disagreement is treated as noise and shrinks the score toward zero.

The fused score is converted into an expected excess return using the standard
information-ratio scaling of \citet{grinold00}, which places the score on a return
scale rather than leaving it as a dimensionless quantity:
\begin{equation}
\alpha_{i,t} = \mathrm{IC} \cdot \sigma_{i,t} \cdot z_{i,t},
\label{eq:alpha}
\end{equation}
where $\sigma_{i,t}$ is the annualised volatility of asset $i$ from the current
covariance estimate and $\mathrm{IC} = 0.05$ is an assumed information coefficient (IC).

Equation~(\ref{eq:alpha}) is not a cosmetic step. If a raw indicator is passed to a
mean-variance objective directly, the units of the linear term and of the quadratic
risk term are unrelated, and their ratio---not the information in the signal---decides
the solution. An indicator such as on-balance volume takes values of order $10^{8}$ to
$10^{9}$, whereas annualised variances are of order $10^{-2}$; the linear term then
dominates by ten orders of magnitude and the optimiser degenerates to a corner
solution, holding a single asset regardless of the risk model. Section~\ref{sec:alloc}
documents this failure mode explicitly, because it is the source of a behaviour visible
in the originally submitted results.

\subsection{Portfolio Construction}

Allocation proceeds in two steps, which separate two decisions that are often conflated:
which assets to hold relative to one another, and how much total risk to carry.

\paragraph{Step 1: relative allocation.} The risky sleeve solves a quadratic program
with a linear transaction-cost penalty,
\begin{align}
\max_{w} \quad & \alpha_t^{\top} w - \lambda \, w^{\top} \Sigma_t w
- \kappa \left\| w - w_{t-1} \right\|_1 \\
\text{s.t.} \quad & w \geq 0, \qquad w \leq w_{\max}, \qquad \mathbf{1}^{\top} w = 1,
\end{align}
where $\Sigma_t$ is the annualised covariance matrix, $\lambda = 2$ is the risk
aversion coefficient, $w_{\max} = 0.25$ caps any single position at one quarter of the
sleeve, and $w_{t-1}$ is the weight vector carried in from the previous rebalance. The
turnover penalty $\kappa = c \cdot 252 / h$ annualises a one-way cost of $c = 10$ basis
points over a rebalancing interval of $h = 21$ sessions, so that the cost of trading is
expressed on the same annual basis as the return and risk terms.

The equality constraint $\mathbf{1}^{\top} w = 1$ is what makes this step purely
relative. It fixes the sleeve to be fully invested, so the optimiser can only express a
view about composition, never about how much risk to take. That second decision is made
explicitly in step~2 instead of emerging as a by-product of the constraint set.

$\Sigma_t$ is estimated by Ledoit--Wolf shrinkage \citep{ledoit04} on the trailing 252
daily returns and annualised by a factor of 252. Shrinkage is used because a sample
covariance matrix estimated from 252 observations on ten assets is poorly conditioned,
and a mean-variance optimiser will concentrate weight precisely on the directions in
which the estimate is least reliable. The problem is convex and is solved with CVXPY
using the Clarabel interior-point solver.

\paragraph{Step 2: regime risk budget.} The sleeve is then scaled by a factor that
depends only on the regime:
\begin{equation}
w^{\ast}_t = k_{r_t} \, \tilde{w}_t, \qquad
k = \{\,\text{calm}: 1.00, \;\; \text{transitional}: 0.75, \;\; \text{stress}: 0.50\,\},
\label{eq:budget}
\end{equation}
with the residual $1 - k_{r_t}$ held in cash, which is assumed to earn zero. Scaling is
one-sided: the framework de-risks in disturbed states but never levers up in calm ones,
so any performance difference is attributable to risk reduction rather than to leverage.

The budget in Equation~(\ref{eq:budget}) is expressed as a fraction of the sleeve's own
risk rather than as an absolute volatility target, and this choice is consequential.
An absolute target de-risks only when the sleeve's estimated volatility exceeds the
target level, so its effect depends on the volatility of the universe it is applied to.
Targets calibrated for a concentrated equity portfolio---of the order of 10 to 20 per
cent annualised---are simply never reached by a diversified ten-asset sleeve whose own
realised volatility is close to 7~per~cent. Under those targets the constraint binds in
one rebalance out of 192 in this study, and the regime layer is inert: it is present in
the specification, consumes computation, and does nothing. A relative budget is
scale-free by construction and remains active in any universe.
Section~\ref{sec:ablation} reports both variants and the difference between them.

\subsection{Walk-Forward Protocol}

The framework is rebalanced every 21 sessions. Between rebalances positions are held
and allowed to drift with realised returns, so that reported turnover reflects trading
actually undertaken rather than a notional daily reweighting. One-way transaction costs
of 10 basis points are charged on the traded fraction $\|w^{\ast}_t - w_{t-1}\|_1$ on
the rebalance session; Section~\ref{sec:cost} varies this assumption from 0 to 25 basis
points.

Every estimated quantity uses only data preceding the decision it informs: the
covariance matrix from the trailing 252 sessions before the rebalance date, the mixture
model from an expanding window strictly earlier than that date, and all features from
causal rolling windows. No parameter is tuned on the evaluation period; the
configuration reported as primary is the one specified above, and
Section~\ref{sec:robust} reports the full grid of alternatives rather than the best
cell of it.

\subsection{Evaluation Design}

Performance is summarised by compound annual growth rate (CAGR), annualised
volatility, Sharpe ratio computed at a zero risk-free rate, maximum drawdown (Max DD)
and Calmar ratio. Three
benchmarks are used. The first is an equal-weighted portfolio of the same ten assets,
rebalanced daily and charged no transaction costs, which is deliberately a generous
comparator. The second is an inverse-volatility portfolio using trailing 60-session
standard deviations, also costless, which isolates whether the framework adds anything
beyond elementary risk weighting. The third rescales the equal-weighted portfolio to
the strategy's realised volatility, which removes the mechanical advantage a
lower-risk portfolio enjoys in ratio comparisons.

Statistical inference on daily excess returns uses a paired $t$-test, a Newey--West
$t$-statistic with ten lags to accommodate autocorrelation and heteroskedasticity, and
a stationary block bootstrap with 21-session blocks and 2{,}000 replications for
interval estimates of the Sharpe ratio and of the Sharpe difference. Component
contributions are isolated by ablation: removing the regime layer while holding average
exposure fixed, removing it while holding full exposure, and removing the signal layer
entirely. Robustness is assessed over a grid of random seeds, covariance estimation
windows and rebalancing frequencies, with every cell reported.

\section{Results}
\label{sec:results}

All results in this section come from the single specification set out in Section~3.
No parameter was selected by reference to the evaluation period, and every
specification examined is reported, including those that are unfavourable to the
framework.

\subsection{Headline Performance}

Table~\ref{tab:main} reports the primary comparison over the 4{,}024 trading sessions from
7~January 2008 to 29~December 2023.

\begin{table}[!ht]
\centering
\caption{Primary results, ten-asset universe, January 2008 to December 2023. Sharpe
ratios use a zero risk-free rate and are computed as compound annual growth divided by
annualised volatility. Benchmarks are rebalanced daily and charged no transaction costs.}
\label{tab:main}
\small
\begin{tabular}{p{4.4cm}rrrrr}
\hline
& \textbf{CAGR} & \textbf{Volatility} & \textbf{Sharpe} & \textbf{Max DD} & \textbf{Calmar} \\
\hline
Regime-conditional strategy & 4.59\% & 6.22\% & 0.738 & $-18.99$\% & 0.242 \\
Equal-weight benchmark & 7.33\% & 13.80\% & 0.532 & $-36.08$\% & 0.203 \\
Inverse-volatility benchmark & 6.74\% & 8.82\% & 0.764 & $-24.09$\% & 0.280 \\
Equal-weight at strategy volatility & 3.48\% & 6.22\% & 0.561 & $-17.67$\% & 0.197 \\
\hline
\end{tabular}
\end{table}

Three readings of Table~\ref{tab:main} are warranted, and they do not all favour the framework.

Against the equal-weighted benchmark the framework improves the Sharpe ratio from 0.532
to 0.738, a gain of 39 per cent, and reduces maximum drawdown from $-36.08$ per cent to
$-18.99$ per cent, a reduction of 47 per cent. It does so while earning a lower
compound return, 4.59 per cent against 7.33 per cent per annum. The framework is not a
return-enhancing strategy and is not presented as one; it is a risk-reduction strategy
whose return sacrifice is smaller than its risk reduction.

Against the volatility-matched equal-weighted benchmark, which removes the mechanical
advantage a lower-risk portfolio enjoys in any ratio comparison, the framework retains
a Sharpe advantage of 0.738 against 0.561, but loses on maximum drawdown, $-18.99$ per
cent against $-17.67$ per cent. Deleveraging a diversified benchmark is therefore a
competitive way to reduce worst-case loss, and the framework's drawdown advantage over
the unlevered benchmark should not be read as evidence that it dominates every
risk-reduced alternative.

Against the inverse-volatility benchmark the framework does not win at all: 0.738
against 0.764 on Sharpe and 0.242 against 0.280 on Calmar, although it does hold a
smaller maximum drawdown, $-18.99$ per cent against $-24.09$ per cent. Elementary risk
weighting is a strong comparator on a diversified multi-asset universe, and this is
reported rather than omitted.

Figure~2 shows the cumulative growth of one unit of capital and Figure~3 the drawdown
profile of each series.

\begin{figure*}[t]
\centering
\includegraphics[width=\textwidth]{images/fig_cumret.png}
\caption{Cumulative growth of one unit of capital, logarithmic scale}
\label{fig:cumret}
\end{figure*}

\begin{figure*}[t]
\centering
\includegraphics[width=\textwidth]{images/fig_drawdown.png}
\caption{Drawdown from running peak}
\label{fig:drawdown}
\end{figure*}

The two figures make the shape of the result visible. The framework's equity curve
tracks below the benchmark through extended advances---most conspicuously 2013, 2017,
2021 and 2023---and separates upward during 2008, 2011, 2020 and 2022. The drawdown
plot shows the same asymmetry more directly: the benchmark reaches $-36$ per cent in
2008 and $-25$ per cent in 2022, while the framework's deepest excursions are
approximately half as large. Recovery is correspondingly faster, since a shallower loss
requires a smaller subsequent gain to restore the previous peak.

\subsection{Robustness Across Specifications}
\label{sec:robust}

The primary configuration fixes a random seed, a covariance estimation window and a
rebalancing interval. Table~\ref{tab:sweep} reports every combination of three seeds, two estimation
windows and three rebalancing intervals, so that the reported result can be located
within the distribution of specifications rather than presented in isolation.

\begin{table}[!ht]
\centering
\caption{Specification sweep. All eighteen cells reported; the last two columns record
whether the cell beats the equal-weight benchmark on each measure. That benchmark attains
a Sharpe ratio of 0.582 and a maximum drawdown of $-36.1$\% on the common evaluation
window used for this comparison.}
\label{tab:sweep}
\small
\begin{tabular}{rrrrrcc}
\hline
\textbf{Seed} & \textbf{Cov. window} & \textbf{Rebalance} & \textbf{Sharpe} & \textbf{Max DD} & \textbf{Sharpe win} & \textbf{DD win} \\
\hline
0 & 126 & 10 & 0.736 & $-20.7$\% & yes & yes \\
0 & 126 & 21 & 0.830 & $-19.6$\% & yes & yes \\
0 & 126 & 42 & 0.802 & $-19.7$\% & yes & yes \\
0 & 252 & 10 & 0.685 & $-20.8$\% & yes & yes \\
0 & 252 & 21 & 0.775 & $-19.0$\% & yes & yes \\
0 & 252 & 42 & 0.776 & $-20.4$\% & yes & yes \\
7 & 126 & 10 & 0.750 & $-20.6$\% & yes & yes \\
7 & 126 & 21 & 0.830 & $-19.6$\% & yes & yes \\
7 & 126 & 42 & 0.788 & $-19.6$\% & yes & yes \\
7 & 252 & 10 & 0.678 & $-21.1$\% & yes & yes \\
7 & 252 & 21 & 0.757 & $-19.0$\% & yes & yes \\
7 & 252 & 42 & 0.763 & $-20.4$\% & yes & yes \\
42 & 126 & 10 & 0.718 & $-20.6$\% & yes & yes \\
42 & 126 & 21 & 0.798 & $-19.6$\% & yes & yes \\
42 & 126 & 42 & 0.765 & $-19.7$\% & yes & yes \\
42 & 252 & 10 & 0.672 & $-21.1$\% & yes & yes \\
42 & 252 & 21 & 0.753 & $-19.0$\% & yes & yes \\
42 & 252 & 42 & 0.721 & $-20.4$\% & yes & yes \\
\hline
\end{tabular}
\end{table}

Every one of the eighteen specifications beats the equal-weighted benchmark on both the
Sharpe ratio and maximum drawdown. The Sharpe ratio ranges from 0.672 to 0.830 and the
maximum drawdown from $-21.1$ to $-19.0$ per cent, so the primary configuration at
0.738 sits near the middle of its own distribution rather than at the top of it.
Sensitivity to the random seed is small, which indicates that the mixture fit is stable;
sensitivity to the rebalancing interval is larger, with ten-session rebalancing
consistently weakest because it incurs more turnover for no additional information.

The same sweep is unfavourable on one comparison. In none of the eighteen cells does
the framework beat the volatility-matched equal-weighted benchmark on maximum drawdown.
This is a consistent finding rather than an artefact of the primary configuration, and
it bounds the drawdown claim precisely: the framework reduces drawdown relative to an
unlevered diversified benchmark, not relative to every portfolio held at the same
volatility.

\subsection{Behaviour in Stress Episodes}

Table~\ref{tab:episodes} reports cumulative returns over four episodes chosen before the analysis on the
basis of documented market history rather than by inspection of the results.

\begin{table}[!ht]
\centering
\caption{Cumulative return over stress episodes. Differences are stated in percentage
points (pp).}
\label{tab:episodes}
\small
\begin{tabular}{p{3.6cm}rrrr}
\hline
\textbf{Episode} & \textbf{Sessions} & \textbf{Strategy} & \textbf{Equal-weight} & \textbf{Difference} \\
\hline
Global financial crisis, Sep 2008--Mar 2009 & 130 & $-7.05$\% & $-29.78$\% & $+22.74$ pp \\
Euro area crisis, Jul--Oct 2011 & 65 & $+3.41$\% & $-7.24$\% & $+10.66$ pp \\
COVID-19 crash, Feb--Mar 2020 & 24 & $-6.29$\% & $-21.28$\% & $+14.99$ pp \\
Inflation bear market, Jan--Oct 2022 & 196 & $-16.75$\% & $-24.49$\% & $+7.74$ pp \\
\hline
\end{tabular}
\end{table}

The framework outperforms in all four episodes, by between 7.7 and 22.7 percentage
points. The mechanism is visible in the exposure series: the mixture model classifies
these periods into the transitional or stress state, and the risk budget reduces
invested exposure accordingly. The margin is widest in 2008, the longest and deepest of
the four, and narrowest in 2022, which is consistent with the design---a gradual,
correlated decline across both equities and bonds offers less scope for a de-risking
rule than an abrupt volatility shock does.

Complementing the episode analysis, the framework's average return on the worst
5~per~cent of benchmark sessions ($n = 201$) is $-0.53$ per cent against $-2.07$ per
cent for the benchmark, a cushion of 1.53 percentage points per session. Downside
capture is 0.35 and upside capture 0.37, giving a capture ratio of 0.93. The ratio
being close to unity is informative: the framework participates less in both directions
in roughly equal proportion, so its advantage comes predominantly from carrying less
risk in adverse states rather than from any directional skill.

\subsection{Year-by-Year Performance}

Table~\ref{tab:years} reports every calendar year in the sample, with both wins and losses shown.

\begin{table}[!ht]
\centering
\caption{Calendar year performance. Sharpe and maximum drawdown computed within each year.}
\label{tab:years}
\small
\begin{tabular}{rrrrrcc}
\hline
\textbf{Year} & \textbf{Strategy} & \textbf{Equal-weight} & \textbf{Strat. Sharpe} & \textbf{EW Sharpe} & \textbf{Sharpe win} & \textbf{DD win} \\
\hline
2008 & $-2.27$\% & $-16.32$\% & $-0.33$ & $-0.47$ & yes & yes \\
2009 & $+4.36$\% & $+24.91$\% & 0.60 & 1.18 & no & yes \\
2010 & $+13.75$\% & $+18.31$\% & 2.05 & 1.40 & yes & yes \\
2011 & $+10.36$\% & $+5.75$\% & 1.73 & 0.45 & yes & yes \\
2012 & $+8.91$\% & $+13.49$\% & 1.96 & 1.58 & yes & yes \\
2013 & $-0.58$\% & $+6.17$\% & $-0.06$ & 0.72 & no & yes \\
2014 & $+13.03$\% & $+9.97$\% & 2.94 & 1.49 & yes & yes \\
2015 & $-0.48$\% & $-1.56$\% & $-0.05$ & $-0.13$ & yes & yes \\
2016 & $+4.83$\% & $+8.49$\% & 0.94 & 0.97 & no & yes \\
2017 & $+10.93$\% & $+16.59$\% & 2.66 & 3.15 & no & yes \\
2018 & $-3.51$\% & $-5.26$\% & $-0.78$ & $-0.52$ & no & yes \\
2019 & $+15.72$\% & $+22.55$\% & 2.84 & 3.03 & no & yes \\
2020 & $+9.64$\% & $+19.32$\% & 1.23 & 0.96 & yes & yes \\
2021 & $-0.80$\% & $+10.15$\% & $-0.09$ & 1.13 & no & no \\
2022 & $-12.98$\% & $-19.30$\% & $-1.55$ & $-1.23$ & no & yes \\
2023 & $+6.83$\% & $+16.53$\% & 0.90 & 1.52 & no & yes \\
\hline
\end{tabular}
\end{table}

The framework wins on Sharpe ratio in 7 of 16 years and on drawdown in 15 of 16. This
distribution is the honest summary of the strategy: it loses on risk-adjusted return in
most individual years, particularly strong advancing years such as 2009, 2017, 2021 and
2023, and recovers the deficit in the minority of years when the benchmark suffers. Two
years merit comment. In 2021 the framework loses on both measures, the only such year,
because the mixture model spent much of that year in the transitional state during what
turned out to be a sustained advance. In 2018 the framework loses on Sharpe despite a
smaller loss in absolute terms, an arithmetic consequence of a smaller negative return
divided by a much smaller volatility.

Figure~4 shows the rolling one-year Sharpe ratio of both series. The framework's rolling
Sharpe is positive in 78.6 per cent of windows against 83.3 per cent for the benchmark,
and exceeds the benchmark in 48.7 per cent of windows. Its advantage is therefore
episodic rather than persistent, which is what Table~\ref{tab:years} leads one to expect.

\begin{figure*}[t]
\centering
\includegraphics[width=\textwidth]{images/fig_rollsharpe.png}
\caption{Rolling 252-session Sharpe ratio}
\label{fig:rollsharpe}
\end{figure*}

\subsection{Component Ablation}
\label{sec:ablation}

The framework has two active layers, signal fusion and the regime risk budget.
Table~\ref{tab:ablation} removes each in turn. The comparison requires care, because scaling a portfolio
toward cash changes its return and its volatility together; comparing a de-risked
strategy against a fully invested variant therefore confounds the timing of risk with
its level. The first ablation holds average invested exposure fixed at the framework's
own 91.9 per cent and applies it constantly, so that only the timing differs.

\begin{table}[!ht]
\centering
\caption{Component ablation. The matched-exposure control holds the same average
invested exposure as the full framework but applies it without regard to the regime.}
\label{tab:ablation}
\small
\begin{tabular}{p{4.8cm}rrrrr}
\hline
\textbf{Configuration} & \textbf{CAGR} & \textbf{Sharpe} & \textbf{Max DD} & \textbf{Exposure} & \textbf{$p$} \\
\hline
Full framework & 4.59\% & 0.738 & $-18.99$\% & 91.9\% & --- \\
No regime timing, same average exposure & 4.58\% & 0.675 & $-22.25$\% & 91.9\% & 0.960 \\
No regime timing, full exposure & 5.79\% & 0.786 & $-22.53$\% & 100.0\% & 0.037 \\
No signal layer & 4.81\% & 0.784 & $-18.87$\% & 91.9\% & 0.313 \\
Equal-weight benchmark & 7.33\% & 0.532 & $-36.08$\% & 100.0\% & --- \\
\hline
\end{tabular}
\end{table}

The first comparison is the one the framework is built on, and it succeeds. Holding
average exposure fixed, allocating that exposure according to the regime rather than
uniformly raises the Sharpe ratio from 0.675 to 0.738 and improves maximum drawdown
from $-22.25$ to $-18.99$ per cent. The difference in mean return is essentially zero
($p = 0.960$), which is the expected signature of a pure risk-timing effect: the same
average risk is carried, but it is carried at better moments.

The second comparison bounds the cost of that timing. Against a fully invested variant,
the framework gives up approximately 1.2 percentage points of annual return, and the difference is
statistically significant ($p = 0.037$). De-risking is not free, and its Sharpe ratio
under this comparison is lower, 0.738 against 0.786. What the framework buys with that
return is 3.54 percentage points of maximum drawdown. Whether that exchange is
worthwhile is a preference, not a fact, and the paper reports both sides of it.

The third comparison is a null, and it concerns the layer the original framework was
named for. Removing the signal fusion layer entirely---setting expected returns to
zero so that the optimiser becomes a constrained minimum-variance problem---changes the
Sharpe ratio from 0.738 to 0.784 and the maximum drawdown from $-18.99$ to $-18.87$ per
cent. Both movements favour removal, and neither is statistically significant
($p = 0.313$). The signal fusion layer contributes nothing detectable to portfolio
outcomes in this universe, and its point estimate is mildly adverse.
Section~\ref{sec:ic} establishes why.

\subsection{Form of the Regime Risk Budget}

Table~\ref{tab:budget} varies the risk budget itself, including the absolute-target formulation used
in the framework as originally specified.

\begin{table}[!ht]
\centering
\caption{Risk budget specifications. Relative budgets are stated as the fraction of the
sleeve's own volatility retained in the calm, transitional and stress states.}
\label{tab:budget}
\small
\begin{tabular}{p{4.6cm}rrrr}
\hline
\textbf{Budget} & \textbf{CAGR} & \textbf{Sharpe} & \textbf{Max DD} & \textbf{Exposure} \\
\hline
Relative 1.00 / 0.75 / 0.50 (primary) & 4.59\% & 0.738 & $-18.99$\% & 91.9\% \\
Relative 1.00 / 0.80 / 0.60 & 4.74\% & 0.737 & $-19.94$\% & 93.5\% \\
Relative 1.00 / 0.60 / 0.30 & 4.23\% & 0.736 & $-16.62$\% & 87.4\% \\
Relative 1.00 / 0.50 / 0.25 & 4.07\% & 0.736 & $-15.71$\% & 84.6\% \\
Absolute targets, 20 / 15 / 10\% annualised & 5.81\% & 0.786 & $-22.86$\% & 100.0\% \\
\hline
\end{tabular}
\end{table}

Two findings follow. First, the Sharpe ratio is almost invariant across the four
relative budgets, 0.736 to 0.738, while maximum drawdown improves monotonically as the
budget tightens, from $-19.94$ to $-15.71$ per cent. The result is therefore not an
artefact of the particular multipliers chosen; the mildest of the four is reported as
primary precisely because it is the least aggressive.

Second, the absolute-target formulation is inert. Its mean invested exposure is 99.99
per cent, and the constraint binds in one rebalance out of 192. Targets of 10 to 20 per
cent annualised volatility are never reached by a sleeve whose own realised volatility
is 6.2 per cent, so under that formulation the regime layer is present in the
specification but has no effect on the portfolio. Its apparently superior Sharpe ratio
of 0.786 is simply the fully invested portfolio of Table~\ref{tab:ablation}, reached by a different
route. This is the practical reason the relative formulation is adopted: an absolute
volatility target is not scale-free, and a target calibrated for one universe silently
disables the mechanism in another.

\subsection{Predictive Content of the Signals}
\label{sec:ic}

Table~\ref{tab:ic} measures the predictive content of the two composites directly, as the
cross-sectional Spearman information coefficient between the score on session $t$ and
the return on session $t+1$, with Newey--West $t$-statistics.

\begin{table}[!ht]
\centering
\caption{Information coefficients against next-session cross-sectional returns.
Newey--West $t$-statistics with ten lags in parentheses.}
\label{tab:ic}
\small
\begin{tabular}{p{3.8cm}rrr}
\hline
& \textbf{Technical} & \textbf{Market state} & \textbf{Fused} \\
\hline
All sessions ($n=4{,}004$) & $-0.036$ ($-2.37$) & $-0.031$ ($-2.60$) & $-0.039$ ($-2.72$) \\
Calm regime ($n=2{,}835$) & $-0.047$ ($-2.58$) & $-0.040$ ($-2.82$) & --- \\
Transitional regime ($n=1{,}043$) & $-0.026$ ($-0.98$) & $-0.016$ ($-0.71$) & --- \\
Stress regime ($n=126$) & $+0.132$ ($+2.13$) & $+0.045$ ($+0.75$) & --- \\
\hline
\end{tabular}
\end{table}

The information coefficients are negative and statistically significant over the full
sample for both composites and for their fusion. A composite whose cross-sectional
score is negatively related to next-session returns cannot improve an allocation that
uses it in the stated direction, and the ablation in Table~\ref{tab:ablation} is the portfolio-level
consequence of the numbers in Table~\ref{tab:ic}.

The regime breakdown rejects the fusion hypothesis specifically. Section~3.5 assigns
greater weight to the market-state composite in the stress regime on the hypothesis
that it becomes relatively more informative there. The data show the opposite. In the
stress regime it is the technical composite that carries information ($+0.132$,
$t = 2.13$) while the market-state composite does not ($+0.045$, $t = 0.75$); in the
calm regime both are significantly negative. The difference between the two composites
moves in the direction opposite to the hypothesis, and the contrast is significant
($t = -2.69$, $p = 0.008$). The regime-conditional weighting schedule is therefore not
merely unhelpful but pointed the wrong way.

Two qualifications are owed. The stress regime contains only 126 sessions, so its
coefficients are imprecisely estimated and the positive technical reading should not be
over-interpreted. And the negative full-sample coefficients are consistent with
short-horizon cross-sectional reversal in a universe of broad index funds, which is a
known effect and not evidence of an error in construction. The conclusion is bounded
accordingly: these composites, as specified, do not carry usable directional
information in this universe at this horizon.

\subsection{Allocation Behaviour, Turnover and Costs}
\label{sec:alloc}
\label{sec:cost}

Table~\ref{tab:turnover} summarises the portfolios the framework actually holds, and Figures~5 and~6
display them.

\begin{table}[!ht]
\centering
\caption{Allocation characteristics and sensitivity to the assumed transaction cost.
Effective number of assets is the reciprocal of the Herfindahl index of sleeve weights.
Annual turnover is one-way, as a multiple of portfolio value.}
\label{tab:turnover}
\small
\begin{tabular}{p{4.8cm}rrrr}
\hline
\multicolumn{5}{l}{\textit{Allocation characteristics, primary specification}} \\
\hline
Mean largest weight & \multicolumn{4}{r}{0.250} \\
Mean effective number of assets & \multicolumn{4}{r}{5.40} \\
Minimum effective number of assets & \multicolumn{4}{r}{4.60} \\
Mean invested exposure & \multicolumn{4}{r}{91.9\%} \\
Minimum invested exposure & \multicolumn{4}{r}{50.0\%} \\
Annual one-way turnover & \multicolumn{4}{r}{0.90$\times$} \\
\hline
\multicolumn{5}{l}{\textit{Sensitivity to assumed cost}} \\
\hline
\textbf{Assumed cost} & \textbf{CAGR} & \textbf{Sharpe} & \textbf{Max DD} & \textbf{Turnover} \\
\hline
0 basis points & 3.96\% & 0.639 & $-19.83$\% & 4.16$\times$ \\
5 basis points & 4.68\% & 0.761 & $-19.54$\% & 1.10$\times$ \\
10 basis points (primary) & 4.59\% & 0.738 & $-18.99$\% & 0.90$\times$ \\
25 basis points & 4.40\% & 0.706 & $-19.29$\% & 0.87$\times$ \\
\hline
\end{tabular}
\end{table}

\begin{figure*}[t]
\centering
\includegraphics[width=\textwidth]{images/fig_alloc.png}
\caption{Portfolio weights over time, including the cash residual}
\label{fig:alloc}
\end{figure*}

\begin{figure*}[t]
\centering
\includegraphics[width=\textwidth]{images/fig_exposure.png}
\caption{Invested exposure over time, with transitional and stress regimes shaded}
\label{fig:exposure}
\end{figure*}

The allocation is diversified and stable. The largest position sits at the 25 per cent
cap on average, the effective number of assets averages 5.40 and never falls below
4.60, and annual one-way turnover is 0.90 times portfolio value, equivalent to
approximately 7.5 per cent of the portfolio traded at each monthly rebalance. Invested
exposure follows the regime, falling to a minimum of 50 per cent during the stress
state, which is the mechanism Figure~6 displays directly.

The cost sensitivity contains a result worth stating explicitly. Performance is not
monotone in the assumed cost: at zero assumed cost the Sharpe ratio is 0.639, below the
0.738 obtained at 10 basis points, even though no costs are deducted. The reason is
that the assumed cost enters the objective as a turnover penalty as well as entering
the return as a deduction. With the penalty removed, annual turnover rises from
0.90 to 4.16 times portfolio value as the optimiser chases small differences in the
estimated covariance matrix, and the resulting weight instability costs more than the
saved commissions are worth. The turnover penalty is therefore functioning as a
regulariser on the weight vector, and the framework is robust across the plausible cost
range of 5 to 25 basis points, over which the Sharpe ratio varies only between 0.706
and 0.761.

This section also resolves a defect in the framework as originally specified. In that
version the fused score entered the objective without the information-ratio scaling of
Equation~(\ref{eq:alpha}), so the linear term was expressed in the raw units of the
underlying indicators. On-balance volume takes values of order $10^{8}$ to $10^{9}$,
while annualised variances are of order $10^{-2}$. The linear term consequently
dominated the quadratic risk term by roughly ten orders of magnitude, the risk model
became irrelevant to the solution, and the optimiser returned a corner solution holding
a single asset on every rebalance date---an effective number of assets of exactly 1.00,
with one holding persisting for 1{,}364 of 1{,}421 sessions. The scaling in
Equation~(\ref{eq:alpha}) removes the defect by placing the linear term on a return
scale, and Table~\ref{tab:turnover} reports the corrected behaviour. The failure was one of units, not
of the optimiser or of the risk model.

\subsection{Regime Model Diagnostics}
\label{sec:regimefit}

Table~\ref{tab:regimefit} reports the mixture model's order selection statistics and the properties of
the three-state solution.

\begin{table}[!ht]
\centering
\caption{Regime model diagnostics. Order selection is computed on the full feature
history; state characteristics and the transition matrix come from the three-component
solution. Out-of-sample shares are the states actually assigned by the expanding-window
refits used in the backtest.}
\label{tab:regimefit}
\small
\begin{tabular}{p{3.2cm}rrrrr}
\hline
\multicolumn{6}{l}{\textit{Order selection}} \\
\hline
\textbf{Components} & 2 & 3 & 4 & 5 & 6 \\
AIC & 19053 & 17900 & 17719 & 17624 & 17481 \\
BIC & 19123 & 18008 & 17866 & 17808 & 17704 \\
\hline
\multicolumn{6}{l}{\textit{Three-state solution}} \\
\hline
& \textbf{Calm} & \textbf{Transitional} & \textbf{Stress} & & \\
Annualised volatility & 12.00\% & 20.34\% & 49.42\% & & \\
Annualised momentum & $+14.97$\% & $-0.19$\% & $-20.11$\% & & \\
Share of sessions, in sample & 61.1\% & 33.7\% & 5.2\% & & \\
Share of sessions, out of sample & 70.8\% & 26.1\% & 3.1\% & & \\
Mean run length (sessions) & 27.7 & 13.9 & 22.1 & & \\
Self-transition probability & 0.964 & 0.928 & 0.955 & & \\
\hline
\multicolumn{6}{l}{\textit{Realised performance within each out-of-sample state}} \\
\hline
& \textbf{Calm} & \textbf{Transitional} & \textbf{Stress} & & \\
Sessions & 2849 & 1049 & 126 & & \\
Strategy Sharpe & 0.971 & 0.233 & 0.877 & & \\
Equal-weight Sharpe & 0.841 & 0.522 & 0.159 & & \\
\hline
\end{tabular}
\end{table}

The order selection statistics are reported rather than used. Both AIC and BIC continue
to fall over the range examined, from 19053 and 19123 at two components to 17481 and
17704 at six, so neither criterion selects an interior optimum and neither supports the
choice of three components. This behaviour is expected: a mixture fitted to persistent,
heavy-tailed financial features will absorb tail observations into additional
components indefinitely, and the criteria reward it for doing so. The choice of three
is made instead on persistence and interpretability, and Table~\ref{tab:regimefit} supplies the grounds.
The three states are economically distinct---annualised volatility of 12.0, 20.3 and
49.4 per cent, with annualised momentum of $+15.0$, $-0.2$ and $-20.1$ per cent---and
they persist, with self-transition probabilities between 0.928 and 0.964 and mean run
lengths of 14 to 28 sessions. States of that duration can be acted upon at a monthly
rebalancing frequency. Finer partitions produce shorter-lived states that the
rebalancing schedule cannot exploit, and the transition matrix is close to
tridiagonal---the estimated probability of moving directly between the calm and the
stress state, in either direction, is zero---which indicates that the ordering recovered by
sorting components on volatility is economically meaningful rather than arbitrary.

The final panel shows where the framework's advantage arises. In the stress state the
framework attains a Sharpe ratio of 0.877 against 0.159 for the benchmark, and in the
calm state 0.971 against 0.841. It underperforms materially only in the transitional
state, 0.233 against 0.522, which is consistent with the 2021 result in Table~\ref{tab:years}: a
state that is genuinely ambiguous causes the risk budget to reduce exposure in
advances that turn out to continue. Figure~7 shows the identified states over time and
Figure~8 the behaviour of the two composites.

\begin{figure*}[t]
\centering
\includegraphics[width=\textwidth]{images/fig_regimes.png}
\caption{Identified market states, with the underlying volatility and momentum features}
\label{fig:regimes}
\end{figure*}

\begin{figure*}[t]
\centering
\includegraphics[width=\textwidth]{images/fig_signal.png}
\caption{Mean absolute cross-sectional score of each composite, 21-session moving average}
\label{fig:signal}
\end{figure*}

\subsection{Statistical Significance}

Table~\ref{tab:inference} reports inference on the daily excess return series and on the Sharpe
difference.

\begin{table}[!ht]
\centering
\caption{Inference against the equal-weighted benchmark. Bootstrap intervals are from a
stationary block bootstrap with 21-session blocks and 2{,}000 replications.}
\label{tab:inference}
\small
\begin{tabular}{p{6.4cm}r}
\hline
Mean excess return over benchmark, annualised & $-3.35$\% \\
Paired $t$-statistic & $-1.211$ \\
$p$-value & 0.226 \\
Newey--West $t$-statistic, ten lags & $-1.511$ \\
\hline
Strategy Sharpe ratio, 95\% bootstrap interval & $[0.213, \; 1.313]$ \\
Benchmark Sharpe ratio, 95\% bootstrap interval & $[0.049, \; 1.095]$ \\
Sharpe difference, median & $+0.203$ \\
Sharpe difference, 95\% bootstrap interval & $[-0.194, \; 0.558]$ \\
Bootstrap probability that the difference is not positive & 0.150 \\
\hline
\end{tabular}
\end{table}

The results must be stated carefully. The framework's mean return is below the
benchmark's, and that shortfall is not statistically significant ($p = 0.226$). The
Sharpe difference of $+0.207$ has a bootstrap median of $+0.203$ and a 95 per cent
interval of $[-0.194, 0.558]$ that includes zero; the bootstrap probability that the
difference is not positive is 0.150. On a sixteen-year sample, therefore, the Sharpe
improvement is not significant at conventional levels, and the paper does not claim
that it is.

What is robust is not the point estimate but its consistency. The Sharpe advantage
appears in 18 of 18 specifications, the drawdown advantage in 18 of 18 and in 15 of 16
calendar years, and the outperformance in all four stress episodes. Sixteen years of
daily data contain few independent observations of the events that matter for a
drawdown claim---four, on the evidence of Table~\ref{tab:episodes}---and no amount of daily sampling
converts that into statistical power. The appropriate conclusion is that the drawdown
reduction is consistent and mechanically explicable, and that the Sharpe improvement,
while positive under every specification examined, is not established at conventional
significance.

\subsection{Secondary Universe: Mega-Capitalisation Equities}
\label{sec:megacap}

The framework was also run without modification on the ten mega-capitalisation US
equities of the original study over 2019--2023, a universe of 1{,}255 sessions.
Table~\ref{tab:megacap} reports the outcome, which is unfavourable.

\begin{table}[!ht]
\centering
\caption{Secondary universe: ten mega-capitalisation US equities, 2019--2023}
\label{tab:megacap}
\small
\begin{tabular}{p{4.8cm}rrrr}
\hline
& \textbf{CAGR} & \textbf{Volatility} & \textbf{Sharpe} & \textbf{Max DD} \\
\hline
Regime-conditional strategy & 13.42\% & 16.58\% & 0.810 & $-20.45$\% \\
Equal-weight benchmark & 33.62\% & 26.68\% & 1.260 & $-33.63$\% \\
Equal-weight at strategy volatility & 20.75\% & 16.58\% & 1.252 & $-21.84$\% \\
No regime timing, same average exposure & 15.78\% & 18.77\% & 0.841 & $-29.75$\% \\
\hline
\end{tabular}
\end{table}

The framework loses decisively on risk-adjusted return, 0.810 against 1.260, while
retaining its defensive properties: maximum drawdown improves from $-33.63$ to $-20.45$
per cent, the framework wins on drawdown in 5 of 5 calendar years but on Sharpe in only
1 of 5, and it outperforms in both stress episodes available in this window, by 12.50
percentage points during the 2020 crash and 12.60 percentage points during the 2022
bear market.
The regime layer continues to function as designed---at matched average exposure it
improves drawdown from $-29.75$ to $-20.45$ per cent---but the improvement is not worth
its cost here.

The reason is a property of the universe rather than of the framework. Ten
mega-capitalisation US equities are close to a single factor bet: correlations are high,
there is no low-volatility asset to rotate into, and the risky sleeve therefore carries
roughly 17 per cent annualised volatility against 6 per cent in the multi-asset case.
De-risking in such a universe can only move capital to cash, so the framework
systematically forgoes return during an exceptional five-year advance and recovers only
part of it in two drawdowns. This bounds the applicability of the framework directly: it
requires a universe with genuine dispersion in asset volatility, and it should not be
expected to add value in a concentrated, single-sector one. The
originally submitted version of this study used exactly such a universe, which is why
its results are superseded here.

\section{Discussion}
\label{sec:discussion}

\subsection{What the Results Establish}

The framework's contribution is narrower than the one claimed for it at the outset, and
better supported. Three findings survive the evaluation of Section~4.

The first concerns the separation of the allocation decision. Splitting the problem into
relative composition and total risk, rather than allowing a single constraint set to
determine both, makes each part independently testable. It is this separation that
allows the ablation in Table~\ref{tab:ablation} to be constructed at all: without it, removing the regime
layer would change both what is held and how much of it is held, and the resulting
comparison would confound the two.

The second is that regime conditioning earns its place through timing, not through
level. At identical average exposure, allocating risk according to the identified state
rather than uniformly improves the Sharpe ratio from 0.675 to 0.738 and maximum
drawdown from $-22.25$ to $-18.99$ per cent, with no difference in mean return
($p = 0.960$). This is the cleanest evidence in the paper, and it is available only
because the control was constructed at matched risk. The corresponding cost is also
measured: against a fully invested variant the framework forgoes approximately 1.2
percentage points of annual return, significantly so, in exchange for 3.5 percentage
points of maximum drawdown.

The third is that the form of the risk budget matters more than its calibration. The
four relative budgets in Table~\ref{tab:budget} produce Sharpe ratios within 0.002 of one another,
while the absolute-target formulation---identical in intent---is inert, binding in one
rebalance out of 192. An absolute volatility target is not scale-free, and a
specification transferred between universes without re-calibration can silently disable
the mechanism it was written to implement. This is a practical warning that generalises
beyond the present framework.

\subsection{What the Results Do Not Establish}

Two claims that the framework was designed to support are not supported.

The signal fusion layer does not contribute. Its removal improves the Sharpe ratio from
0.738 to 0.784 and leaves maximum drawdown unchanged, and neither movement is
significant. The information coefficients in Table~\ref{tab:ic} explain why: both composites are
significantly negatively related to next-session cross-sectional returns over the full
sample. The regime-conditional weighting schedule fares worse still. It was specified on
the hypothesis that market-state information becomes relatively more informative under
stress; in the data the opposite holds, significantly so ($t = -2.69$, $p = 0.008$).

The improvement in risk-adjusted return is not statistically significant. The Sharpe
difference of $+0.207$ carries a bootstrap interval of $[-0.194, 0.558]$. Sixteen years
of daily returns contain only a handful of independent stress events, and a claim about
behaviour during market stress cannot be established at conventional significance from
four episodes. The consistency of the result across 18 of 18 specifications, 15 of 16
calendar years and 4 of 4 episodes is the stronger evidence, and it is evidence of
robustness rather than of significance. The two should not be conflated, and this paper
does not conflate them.

\subsection{Implications for Practice}

For a practitioner the results suggest a specific and limited use. A regime-conditional
risk budget is a reasonable component of a multi-asset allocation process where the
objective includes drawdown control and where some return may be sacrificed for it. It
should be implemented in relative rather than absolute form, evaluated against a
volatility-matched benchmark rather than an unlevered one, and not expected to help in
a universe without genuine dispersion in asset volatility.

The results also caution against a common design pattern. The framework as originally
specified contained two layers whose contribution had never been isolated, one of which
was inert by construction and the other of which was pointed in the wrong direction.
Neither defect was visible in aggregate performance, and both were found only by
ablation at matched risk. Architectures assembled from plausible components should be
disassembled before they are trusted.

Finally, the turnover result in Table~\ref{tab:turnover} has a practical reading. The assumed
transaction cost enters the objective as a penalty as well as the return as a
deduction, so setting it to zero does not produce the best achievable result---it
produces a noisier weight vector and a lower Sharpe ratio. A modest cost assumption acts
as a regulariser, and there is no reason to model a frictionless world even when
evaluating one.

\subsection{Limitations}

Several limitations bound what can be concluded.

The universe is small. Ten exchange-traded funds are enough to demonstrate the
mechanism but not to establish that it generalises to larger or differently composed
universes, and the position cap of 25 per cent binds on average, which means the
optimiser's freedom is partly determined by the constraint rather than by the risk
model.

The evaluation is a single historical path. Sixteen years is a long sample by the
standards of this literature and a short one for inference about tail events, and the
four stress episodes are not independent draws from a stationary process.

The signals are market-derived proxies rather than textual sentiment. No news, filings
or social media text is used at any point, as stated in Section~3.3. The null reported
in Section~4.7 is therefore a null about these proxies, and it is not evidence that
textual sentiment would fail in the same role. Establishing that would require a
time-stamped news corpus with coverage of the traded instruments, which for broad
exchange-traded funds does not exist at the instrument level.

The regime model is deliberately simple. A three-component Gaussian mixture on two
features is interpretable and stable, but it imposes conditional normality within
states, treats observations as conditionally independent given the state, and---as
Table~\ref{tab:regimefit} shows---is not the specification either information criterion would select.

The risk-free rate is taken as zero throughout, and cash is assumed to earn nothing.
Over a sample containing both the near-zero rates of 2009--2015 and the higher rates of
2023, this understates the return on the framework's cash holding and therefore
understates its performance relative to the fully invested benchmark. The direction of
this bias is against the framework, which is the conservative choice, but it is a
simplification.

\subsection{Future Work}

Four extensions follow directly. The first is textual sentiment, applied to a universe
where an instrument-level news corpus exists, which would convert the proxy null of
Section~4.7 into a genuine test of the fusion hypothesis. The second is a larger and
more heterogeneous universe, including commodities, currencies and non-US fixed income,
which would test whether the volatility-dispersion condition identified in
Section~4.11 is the binding requirement it appears to be. The third is learning the
risk-budget multipliers rather than fixing them, for which reinforcement learning is a
natural candidate given the sequential structure of the problem, subject to the
overfitting concerns documented by \citet{harvey15} and \citet{valls26}. The fourth is
replacing the mixture model with a formulation that represents state persistence
explicitly, such as a hidden Markov or jump model \citep{li25}, which would remove the
conditional-independence assumption noted above.

\section{Conclusion}
\label{sec:conclusion}

This paper set out to close the step between a predictive signal and an executable
portfolio, and to determine which parts of a regime-conditional framework are
responsible for its behaviour. Both objectives are met, though not in the direction
originally anticipated.

The framework, evaluated on ten liquid multi-asset instruments over sixteen years under
a walk-forward protocol with transaction costs, attains a Sharpe ratio of 0.738 against
0.532 for an equal-weighted benchmark and reduces maximum drawdown from $-36.1$ to
$-19.0$ per cent, while earning a lower compound return. It outperforms in all four
major stress episodes in the sample, improves on drawdown in 15 of 16 calendar years,
and beats the benchmark on both measures in all 18 specifications examined. The Sharpe
improvement is consistent but not statistically significant: its bootstrap interval
includes zero.

Ablation locates the source of that behaviour precisely. At identical average exposure,
timing risk by the identified market state improves both the Sharpe ratio and the
maximum drawdown with no difference in mean return, which is the signature of a genuine
risk-timing effect. The form of the risk budget proves more consequential than its
calibration: expressed as a fraction of the portfolio's own volatility the layer is
active and effective, while the absolute volatility targets of the original
specification bind in one rebalance out of 192 and leave it inert.

The signal fusion layer, which gave the framework its name, does not contribute.
Removing it improves the Sharpe ratio, and the information coefficients of both
composites are significantly negative over the full sample. The regime-conditional
weighting schedule is worse than neutral: it was specified on the hypothesis that
market-state information gains importance under stress, and the data reject that
hypothesis in the opposite direction at $p = 0.008$. This null is reported in full,
with the coefficients that produce it, because a component-level negative result is more
useful than an aggregate number that conceals which parts of an architecture are doing
the work.

The study's contribution is therefore a specification and a method of testing it: a
two-step separation of composition from total risk that makes each layer independently
falsifiable, an ablation design that isolates risk timing from risk level by matching
average exposure, and a documented account of two failure modes---an inert absolute risk
budget, and a signal-to-allocation conversion in which mismatched units rather than
information determine the solution. The framework is useful within stated bounds: it
requires a universe with genuine dispersion in asset volatility, it trades return for
drawdown at a measured rate, and it should be evaluated against a volatility-matched
benchmark rather than an unlevered one. Outside those bounds, as the mega-capitalisation
results show, it does not add value.


\section*{Acknowledgements}
\label{sec:declarations}
The authors would like to thank the Fr.~C.~Rodrigues Institute of Technology, Vashi and
Fr.~C.~Rodrigues College of Engineering, Bandra for their support.

\section*{Author contributions}
\textbf{Anandkumar Pardeshi:} Conceptualization, Methodology, Software, Data curation,
Formal analysis, Writing -- original draft.\\
\textbf{Sujata Deshmukh:} Supervision, Validation, Writing -- review \& editing,
Research coordination.

\section*{Conflict of interest}
The authors declare no conflict of interest.

\section*{Data and code availability}
All market data used in this study are publicly available daily price and volume series
obtained from Yahoo Finance for the instruments listed in Section~3.1. The analysis code
implementing the framework, the specification sweep and the ablations is available from
the corresponding author on reasonable request.

\section*{AI tool usage declaration}
During the preparation of this work, the authors used [AUTHORS: NAME THE TOOL(S)
ACTUALLY USED] in order to improve the readability and language of the manuscript. After
using this tool, the authors reviewed and edited the content as needed and take full
responsibility for the content of the publication. No part of the analysis, the
generation of results, or the interpretation of findings was delegated to such tools.

\label{sec:refs}
\begin{thebibliography}{46}

\bibitem[\protect\citeauthoryear{Agal et al.}{2025}]{agal25}
Agal, S., Raulji, K. and Odedra, N.D. (2025) 'A machine learning approach to risk based asset allocation in portfolio optimization', {\it Scientific Reports}, Vol. 15, No. 1, doi: 10.1038/s41598-025-26337-x.

\bibitem[\protect\citeauthoryear{Anbaee Farimani et al.}{2024}]{farimani24}
Anbaee Farimani, S., Jahan, M.V. and Milani Fard, A. (2024) 'An Adaptive Multimodal Learning Model for Financial Market Price Prediction', {\it IEEE Access}, Vol. 12, pp. 121846--121863, doi: 10.1109/ACCESS.2024.3441029.

\bibitem[\protect\citeauthoryear{Ang and Bekaert}{2002}]{angbekaert02}
Ang, A. and Bekaert, G. (2002) 'International Asset Allocation With Regime Shifts', {\it Review of Financial Studies}, Vol. 15, No. 4, pp. 1137--1187, doi: 10.1093/rfs/15.4.1137.

\bibitem[\protect\citeauthoryear{Ashrafzadeh et al.}{2025}]{ashrafzadeh25}
Ashrafzadeh, M., Sadrani, M. and Zolfani, S.H. (2025) 'Deep learning and machine learning models for portfolio optimization: Enhancing return prediction with stock clustering', {\it Results in Engineering}, Vol. 27, pp. 106263, doi: 10.1016/j.rineng.2025.106263.

\bibitem[\protect\citeauthoryear{Asimit et al.}{2025}]{asimit25}
Asimit, V., Chong, W.F., Tunaru, R. and Zhou, F. (2025) 'Portfolio selection and risk sharing via risk budgeting', {\it Insurance: Mathematics and Economics}, Vol. 125, pp. 103139, doi: 10.1016/j.insmatheco.2025.103139.

\bibitem[\protect\citeauthoryear{Barroso and Santa-Clara}{2015}]{barroso15}
Barroso, P. and Santa-Clara, P. (2015) 'Momentum has its moments', {\it Journal of Financial Economics}, Vol. 116, No. 1, pp. 111--120, doi: 10.1016/j.jfineco.2014.11.010.

\bibitem[\protect\citeauthoryear{Chou and Nguyen}{2018}]{chou18}
Chou, J.S. and Nguyen, T.K. (2018) 'Forward Forecast of Stock Price Using Sliding-Window Metaheuristic-Optimized Machine-Learning Regression', {\it IEEE Transactions on Industrial Informatics}, Vol. 14, No. 7, pp. 3132--3142, doi: 10.1109/TII.2018.2794389.

\bibitem[\protect\citeauthoryear{Chuang et al.}{2025}]{chuang25}
Chuang, H.M., He, H.C. and Hu, M.C. (2025) 'Chinese Financial News Analysis for Sentiment and Stock Prediction: A Comparative Framework with Language Models', {\it Big Data and Cognitive Computing}, Vol. 9, No. 10, pp. 263, doi: 10.3390/bdcc9100263.

\bibitem[\protect\citeauthoryear{DeMiguel et al.}{2009}]{demiguel09}
DeMiguel, V., Garlappi, L. and Uppal, R. (2009) 'Optimal Versus Naive Diversification: How Inefficient is the $1/N$ Portfolio Strategy?', {\it Review of Financial Studies}, Vol. 22, No. 5, pp. 1915--1953, doi: 10.1093/rfs/hhm075.

\bibitem[\protect\citeauthoryear{Grinold and Kahn}{2000}]{grinold00}
Grinold, R.C. and Kahn, R.N. (2000) {\it Active Portfolio Management: A Quantitative Approach for Producing Superior Returns and Controlling Risk}, 2nd ed., McGraw-Hill, New York, NY.

\bibitem[\protect\citeauthoryear{Gu et al.}{2020}]{gu20}
Gu, S., Kelly, B. and Xiu, D. (2020) 'Empirical Asset Pricing via Machine Learning', {\it The Review of Financial Studies}, Vol. 33, No. 5, pp. 2223--2273, doi: 10.1093/rfs/hhaa009.

\bibitem[\protect\citeauthoryear{Ha et al.}{2026}]{ha26}
Ha, Y.S., Kang, J., Kim, J. and Chun, D. (2026) 'Machine learning-based portfolio optimization: comparative analysis with the all-weather portfolio strategy', {\it Financial Innovation}, Vol. 12, No. 1, doi: 10.1186/s40854-026-00927-8.

\bibitem[\protect\citeauthoryear{Hamilton}{1989}]{hamilton89}
Hamilton, J.D. (1989) 'A New Approach to the Economic Analysis of Nonstationary Time Series and the Business Cycle', {\it Econometrica}, Vol. 57, No. 2, pp. 357--384, doi: 10.2307/1912559.

\bibitem[\protect\citeauthoryear{Harvey and Liu}{2015}]{harvey15}
Harvey, C.R. and Liu, Y. (2015) 'Backtesting', {\it The Journal of Portfolio Management}, Vol. 42, No. 1, pp. 13--28, doi: 10.3905/jpm.2015.42.1.013.

\bibitem[\protect\citeauthoryear{Hood and Raughtigan}{2025}]{hood25}
Hood, B.T. and Raughtigan, C. (2025) 'Volatility Targeting Is Trendy: How Trend Following Explains Alpha in Volatility-Managed Strategies', {\it The Journal of Portfolio Management}, Vol. 52, No. 1, pp. 100--121, doi: 10.3905/jpm.2025.1.764.

\bibitem[\protect\citeauthoryear{Huang et al.}{2020}]{huang20}
Huang, Q., Yang, J., Feng, X., Liew, A.W.C. and Li, X. (2020) 'Automated Trading Point Forecasting Based on Bicluster Mining and Fuzzy Inference', {\it IEEE Transactions on Fuzzy Systems}, Vol. 28, No. 2, pp. 259--272, doi: 10.1109/TFUZZ.2019.2904920.

\bibitem[\protect\citeauthoryear{Jiang et al.}{2025}]{jiang25}
Jiang, Y., Olmo, J. and Atwi, M. (2025) 'High-dimensional multi-period portfolio allocation using deep reinforcement learning', {\it International Review of Economics \& Finance}, Vol. 98, pp. 103996, doi: 10.1016/j.iref.2025.103996.

\bibitem[\protect\citeauthoryear{Kehinde et al.}{2025}]{kehinde25}
Kehinde, T.O., Chung, S.H. and Olanrewaju, O.A. (2025) 'Inverse DEA for Portfolio Volatility Targeting: Industry Evidence from Taiwan Stock Exchange', {\it International Journal of Financial Studies}, Vol. 13, No. 4, pp. 192, doi: 10.3390/ijfs13040192.

\bibitem[\protect\citeauthoryear{Khodamoradi et al.}{2023}]{khodamoradi23}
Khodamoradi, T., Salahi, M. and Najafi, A.R. (2023) 'Multi-intervals robust mean-conditional value-at-risk portfolio optimisation with conditional scenario reduction technique', {\it International Journal of Applied Decision Sciences}, Vol. 16, No. 2, pp. 237, doi: 10.1504/ijads.2023.129475.

\bibitem[\protect\citeauthoryear{Kritzman et al.}{2012}]{kritzman12}
Kritzman, M., Page, S. and Turkington, D. (2012) 'Regime Shifts: Implications for Dynamic Strategies', {\it Financial Analysts Journal}, Vol. 68, No. 3, pp. 22--39, doi: 10.2469/faj.v68.n3.3.

\bibitem[\protect\citeauthoryear{Ledoit and Wolf}{2004}]{ledoit04}
Ledoit, O. and Wolf, M. (2004) 'A well-conditioned estimator for large-dimensional covariance matrices', {\it Journal of Multivariate Analysis}, Vol. 88, No. 2, pp. 365--411, doi: 10.1016/s0047-259x(03)00096-4.

\bibitem[\protect\citeauthoryear{Ledoit and Wolf}{2008}]{ledoit08}
Ledoit, O. and Wolf, M. (2008) 'Robust performance hypothesis testing with the Sharpe ratio', {\it Journal of Empirical Finance}, Vol. 15, No. 5, pp. 850--859, doi: 10.1016/j.jempfin.2008.03.002.

\bibitem[\protect\citeauthoryear{Lee et al.}{2007}]{lee07}
Lee, J.W., Park, J., O, J., Lee, J. and Hong, E. (2007) 'A Multiagent Approach to $Q$-Learning for Daily Stock Trading', {\it IEEE Transactions on Systems, Man, and Cybernetics - Part A: Systems and Humans}, Vol. 37, No. 6, pp. 864--877, doi: 10.1109/TSMCA.2007.904825.

\bibitem[\protect\citeauthoryear{Li et al.}{2025}]{li25}
Li, X., Chen, J., Tao, X. and Ji, Y. (2025) 'Regime-Switching Asset Allocation Using a Framework Combing a Jump Model and Model Predictive Control', {\it Mathematics}, Vol. 13, No. 17, pp. 2837, doi: 10.3390/math13172837.

\bibitem[\protect\citeauthoryear{Loughran and Mcdonald}{2011}]{loughran11}
Loughran, T. and Mcdonald, B. (2011) 'When Is a Liability Not a Liability? Textual Analysis, Dictionaries, and 10‐Ks', {\it The Journal of Finance}, Vol. 66, No. 1, pp. 35--65, doi: 10.1111/j.1540-6261.2010.01625.x.

\bibitem[\protect\citeauthoryear{Lu and Tian}{2025}]{lu25}
Lu, W. and Tian, Z. (2025) 'Nonlinear Shrinkage Estimation of Higher-Order Moments for Portfolio Optimization Under Uncertainty in Complex Financial Systems', {\it Entropy}, Vol. 27, No. 10, pp. 1083, doi: 10.3390/e27101083.

\bibitem[\protect\citeauthoryear{Luo and Mulvey}{2026}]{luo26}
Luo, Y. and Mulvey, J.M. (2026) 'Regime-Aware Asset Allocation with Dual-Regime Signals and Regime-Dependent Asset Selection', {\it The Journal of Financial Data Science}, Vol. 8, No. 3, pp. 104--134, doi: 10.3905/jfds.2026.013.

\bibitem[\protect\citeauthoryear{López de Prado}{2016}]{prado16}
López de Prado, M. (2016) 'Building Diversified Portfolios that Outperform Out of Sample', {\it The Journal of Portfolio Management}, Vol. 42, No. 4, pp. 59--69, doi: 10.3905/jpm.2016.42.4.059.

\bibitem[\protect\citeauthoryear{Markowitz}{1952}]{markowitz52}
Markowitz, H. (1952) 'Portfolio Selection', {\it The Journal of Finance}, Vol. 7, No. 1, pp. 77--91, doi: 10.2307/2975974.

\bibitem[\protect\citeauthoryear{Mazzotta and Baidoo}{2025}]{mazzotta25}
Mazzotta, S. and Baidoo, E. (2025) 'On the fringe of credit visibility: the value of alternative data for assessing the credit risk of subprime underbanked consumers', {\it International Journal of Applied Decision Sciences}, Vol. 18, No. 4, doi: 10.1504/ijads.2025.10066241.

\bibitem[\protect\citeauthoryear{Moreira and Muir}{2017}]{moreira17}
Moreira, A. and Muir, T. (2017) 'Volatility‐Managed Portfolios', {\it The Journal of Finance}, Vol. 72, No. 4, pp. 1611--1644, doi: 10.1111/jofi.12513.

\bibitem[\protect\citeauthoryear{Muhammad et al.}{2026}]{muhammad26}
Muhammad, I., Rospocher, M., Knez, T. and Žitnik, S. (2026) 'Benchmarking large language models for target-based financial sentiment and stock return', {\it International Journal of Data Science and Analytics}, Vol. 22, No. 1, doi: 10.1007/s41060-026-01201-x.

\bibitem[\protect\citeauthoryear{Mun and Kim}{2025}]{mun25}
Mun, Y. and Kim, N. (2025) 'Leveraging Large Language Models for Sentiment Analysis and Investment Strategy Development in Financial Markets', {\it Journal of Theoretical and Applied Electronic Commerce Research}, Vol. 20, No. 2, pp. 77, doi: 10.3390/jtaer20020077.

\bibitem[\protect\citeauthoryear{Newey and West}{1987}]{newey87}
Newey, W.K. and West, K.D. (1987) 'A Simple, Positive Semi-Definite, Heteroskedasticity and Autocorrelation Consistent Covariance Matrix', {\it Econometrica}, Vol. 55, No. 3, pp. 703--708, doi: 10.2307/1913610.

\bibitem[\protect\citeauthoryear{Officioso et al.}{2026}]{officioso26}
Officioso, G., Sperlì, G. and Vignali, A. (2026) 'Financial news sentiment meets market data: A large language model-based approach to stock price prediction', {\it Information Sciences}, Vol. 754, pp. 123711, doi: 10.1016/j.ins.2026.123711.

\bibitem[\protect\citeauthoryear{Politis and Romano}{1994}]{politis94}
Politis, D.N. and Romano, J.P. (1994) 'The Stationary Bootstrap', {\it Journal of the American Statistical Association}, Vol. 89, No. 428, pp. 1303--1313, doi: 10.1080/01621459.1994.10476870.

\bibitem[\protect\citeauthoryear{Raimundo and Bravo}{2026}]{raimundo26}
Raimundo, B. and Bravo, J.M. (2026) 'Forecasting meets Portfolio Theory: A Bibliometric Approach to Decision-Making under Uncertainty', {\it Humanities and Social Sciences Communications}, Vol. 13, No. 1, doi: 10.1057/s41599-026-07252-6.

\bibitem[\protect\citeauthoryear{Saijai and Pankwaen}{2026}]{saijai26}
Saijai, W. and Pankwaen, K. (2026) 'Deep reinforcement learning for portfolio allocation: a comparative study on five Vietnamese equities', {\it Asian Journal of Economics and Banking}, Vol. 10, No. 2, pp. 174--191, doi: 10.1108/ajeb-10-2025-0168.

\bibitem[\protect\citeauthoryear{Sami et al.}{2025}]{sami25}
Sami, H.M., Shah Jalal, M., Faisal Kabir, M. and Nurul Huda, M. (2025) 'IntPort: An Intelligent Portfolio Construction Technique Based on Financial Forecasting by Statistical Average Method', {\it IEEE Access}, Vol. 13, pp. 35355--35375, doi: 10.1109/ACCESS.2025.3543741.

\bibitem[\protect\citeauthoryear{Sheppert}{2026}]{sheppert26}
Sheppert, A.P. (2026) 'The GT-Score: A Robust Objective Function for Reducing Overfitting in Data-Driven Trading Strategies', {\it Journal of Risk and Financial Management}, Vol. 19, No. 1, pp. 60, doi: 10.3390/jrfm19010060.

\bibitem[\protect\citeauthoryear{Singsiri and Yokrattanasak}{2026}]{singsiri26}
Singsiri, P. and Yokrattanasak, J. (2026) 'Machine Learning-Driven Portfolio Optimization Using Money Flow Index-Based Sentiment Signals', {\it International Journal of Financial Studies}, Vol. 14, No. 5, pp. 112, doi: 10.3390/ijfs14050112.

\bibitem[\protect\citeauthoryear{Valls and Steurer}{2026}]{valls26}
Valls, F.B. and Steurer, E. (2026) 'When Strategies Work Too Well: Volatility-Based Investing, Overfitting and the Limits of Empirical Finance', {\it Journal of Financial Risk Management}, Vol. 15, No. 03, pp. 189--211, doi: 10.4236/jfrm.2026.153012.

\bibitem[\protect\citeauthoryear{Wang et al.}{2026}]{wang26}
Wang, H., Xie, D., Duan, Y., Xiong, W. and Chen, D. (2026) 'Optimizing stock market prediction and stock trading strategies with deep learning models enhanced by nonlinear feature identification and robust prediction evaluation', {\it Financial Innovation}, Vol. 12, No. 1, doi: 10.1186/s40854-026-00929-6.

\bibitem[\protect\citeauthoryear{Zhang et al.}{2022}]{zhang22}
Zhang, C., Sjarif, N.N.A. and Ibrahim, R.B. (2022) 'Decision Fusion for Stock Market Prediction: A Systematic Review', {\it IEEE Access}, Vol. 10, pp. 81364--81379, doi: 10.1109/ACCESS.2022.3195942.

\bibitem[\protect\citeauthoryear{Zhang et al.}{2025}]{zhangz25}
Zhang, Z., Choo, W.C. and Arasan, J. (2025) 'Does the optimal model always perform the best A combined approach for interval forecasting', {\it International Journal of Applied Decision Sciences}, Vol. 18, No. 1, pp. 64--83, doi: 10.1504/ijads.2025.143079.

\bibitem[\protect\citeauthoryear{Zong et al.}{2025}]{zong25}
Zong, C., Wan, J., Cascone, L. and Zhou, H. (2025) 'Stock movement prediction with multimodal stable fusion via gated cross-attention mechanism', {\it Complex \& Intelligent Systems}, Vol. 11, No. 9, doi: 10.1007/s40747-025-02023-3.

\end{thebibliography}
\end{document}
